% GENERATED FILE. Edit canonical lessons, then run npm run generate:books.
% canonical-source-hash: fnv1a64:8a2bdd2f7c347634
% canonical-lessons: SA-C08-cintayati, SA-C08-avagacchati, SA-C08-pathati, SA-C08-likhati, SA-S01-letter-ma, SA-S107-letter-ta

\chapter{The Mind and the Palm Leaf}
\label{ch:sa-mind-verbs}
\hlchaptermodality{8}

\begin{chapteropening}
\textbf{By the end of this chapter:} \emph{I can say that someone thinks, understands, reads and writes in Sanskrit; I can name the class that plants -\sk{अय}- inside a verb, build \textquotedblleft{}understands\textquotedblright{} out of a prefix and a verb I already had, say what \sk{लिख्} meant before it meant write, and tell a word worn down by speech (\sk{तद्भव}) from one carried over whole (\sk{तत्सम}).}

Say \sk{लिखति}, give the scratching sense its root carried first, run the chapter's four verbs back as root plus class plus ending — the planted -\sk{अय}- of \sk{चिन्तयति}, the \sk{अव}- prefix and the waking \sk{बुध्} behind understanding, the plain -a- of \sk{पठति} and \sk{लिखति} — sort Bengali lekhā and Hindi cintā into \sk{तद्भव} and \sk{तत्सम}, and name the two word-histories that needed a label rather than a claim.
\end{chapteropening}

\section[cintayati]{\sk{चिन्तयति} (cintayati) — \textquotedblleft{}thinks,\textquotedblright{} and a root that never reached English}

\label{lesson:SA-C08-cintayati}



\begin{quote}
Say \textquotedblleft{}it knows.\textquotedblright{} (\emph{Jānāti}, root \textbf{\sk{ज्ञा}} — the root English \emph{know} is made of.) Here is what you do \emph{before} you know.
\end{quote}

\subsection*{You'll want to know: \sk{चिन्तयति}}
\noindent
\begin{tabularx}{\linewidth}{@{}>{\raggedright\arraybackslash}X>{\raggedright\arraybackslash}X>{\raggedright\arraybackslash}X@{}}
\toprule
\textbf{Sanskrit} & \textbf{sound} & \textbf{meaning} \\
\midrule
\sk{चिन्तयति} & \emph{cintayati} & he, she, or it \textbf{thinks} \\
\sk{चिन्ता} & \emph{cintā} & a \textbf{thought}, a worry \\
\bottomrule
\end{tabularx}

The root is \textbf{\sk{चिन्त्}} (\emph{cint}), \textquotedblleft{}to think, to turn over.\textquotedblright{} Chapter 3 already handed you the noun: \textbf{\sk{न} \sk{चिन्ता}} (\emph{na cintā}), \textquotedblleft{}no worry,\textquotedblright{} is this root saying \textquotedblleft{}think nothing of it.\textquotedblright{}

\begin{sounds}
\textbf{\sk{चि}} is \emph{ci}, said \emph{chi} as in \emph{chip}. Then \textbf{\sk{न्त}}, \textbf{\sk{न}} \emph{na} stacked with \textbf{\sk{त}} \emph{ta} and said as one \emph{nt}. \textbf{\sk{य}} is \emph{ya}, \textbf{\sk{ति}} is \emph{ti}: \emph{chin-ta-ya-ti}.
\end{sounds}

\begin{grammarlens}[title={Grammar Lens: the tenth class plants \sk{अय}}]
Three ways a \sk{गण} class has turned a root into a present stem: add nothing (\emph{as-ti}), add a vowel (\emph{bhava-ti}), plant a syllable inside (\emph{jā-nā-ti}). Here is the fourth.

The \textbf{tenth} class inserts \textbf{-\sk{अय}-} (\emph{-aya-}) between root and ending: \emph{cint} + \emph{aya} + \emph{ti} $\to$ \textbf{\sk{चिन्तयति}}. The marker is loud and always in the same position, so \emph{-aya-} names its class on hearing.
\end{grammarlens}

\begin{cousinweb}
\textbf{\sk{चिन्त्}} goes back to PIE *\emph{kweyt-}, \textquotedblleft{}to notice\textquotedblright{} — and its \textbf{\sk{च}} \emph{c-} is the rounded \emph{k} from the numbers, which Indo-Iranian merged with plain \emph{k} and then palatalised before a front vowel.

Where it went next is the surprise. Baltic and Slavic took \textquotedblleft{}notice\textquotedblright{} to \textquotedblleft{}count\textquotedblright{} and then to \textbf{read}: Russian \emph{čitat} and Lithuanian \emph{skaityti} both mean \textquotedblleft{}to read.\textquotedblright{} Sanskrit made the same root mean \emph{think}.

And English? Nothing descends from *\emph{kweyt-}. The trail west runs out, and saying so beats inventing a cousin.
\end{cousinweb}

\subsection*{Your turn}
Say these aloud:

\begin{itemize}
\raggedright
  \item \textquotedblleft{}cintayati\textquotedblright{} — it thinks
  \item the root and its planted syllable — \emph{cint} … \emph{aya}
  \item the Chapter 3 phrase built on the same root — \emph{na cintā}
\end{itemize}

\begin{tcolorbox}[breakable,colback=teal!4,colframe=teal!35!black,title={Before you move on}]
What does class 10 put between root and ending, and what is \emph{cintayati}'s root? (\textbf{-\sk{अय}-}; \textbf{\sk{चिन्त्}}, the root of \emph{na cintā}.) Why does the word begin with \emph{c} and not \emph{k}? (\textbf{A rounded PIE \emph{k}, merged and then palatalised} — the numbers' rule.) Which verb of knowing came just before this one, and does \emph{cint} have an English cousin? (\textbf{\sk{जानाति}}; \textbf{no} — this root left English nothing.)
\end{tcolorbox}
\section[avagacchati · budhyate]{\sk{अवगच्छति} (avagacchati) — \textquotedblleft{}understands,\textquotedblright{} going down into it}

\label{lesson:SA-C08-avagacchati}



\begin{quote}
Say \textquotedblleft{}it goes,\textquotedblright{} then \textquotedblleft{}it comes.\textquotedblright{} (\emph{Gacchati}, \emph{āgacchati} — one \sk{उपसर्ग} prefix apart.) And say \textquotedblleft{}it thinks.\textquotedblright{} (\emph{Cintayati}.) One more prefix, and thinking turns into understanding.
\end{quote}

\subsection*{You'll want to know: \sk{अवगच्छति}}
\noindent
\begin{tabularx}{\linewidth}{@{}>{\raggedright\arraybackslash}X>{\raggedright\arraybackslash}X>{\raggedright\arraybackslash}X@{}}
\toprule
\textbf{Sanskrit} & \textbf{sound} & \textbf{meaning} \\
\midrule
\sk{गच्छति} & \emph{gacchati} & it \textbf{goes} \\
\sk{अवगच्छति} & \emph{avagacchati} & it \textbf{understands} \\
\bottomrule
\end{tabularx}

\textbf{\sk{अव}} (\emph{ava}) means \textquotedblleft{}down.\textquotedblright{} Put it on \textquotedblleft{}goes\textquotedblright{} and you get \textquotedblleft{}goes down into\textquotedblright{} — and that is Sanskrit's ordinary verb for understanding. No new root, no new stem: one more preverb on a word you already own.

Sanskrit has a second verb for it too, \textbf{\sk{बुध्यते}} (\emph{budhyate}), \textquotedblleft{}wakes, comes to realise.\textquotedblright{} Its ending \textbf{-\sk{ते}} \emph{-te} belongs to Sanskrit's second set of endings, used when the action loops back on the doer — waking is something you do to yourself.

\begin{sounds}
Everything after two syllables is last chapter's word. The new opening is \textbf{\sk{अ}} \emph{a} plus \textbf{\sk{व}} \emph{va}: \emph{a-va-gach-cha-ti}. In \textbf{\sk{बुध्यते}} the middle is \textbf{\sk{ध्य}}, a stack of \textbf{\sk{ध}} \emph{dha} and \textbf{\sk{य}} \emph{ya}, said glued: \emph{bu-dhya-te}.
\end{sounds}

\begin{cousinweb}
\textbf{\sk{गम्}}, under \emph{avagacchati}, is still PIE *\emph{gwem-}, the root of English \textbf{come} and Latin \emph{venīre} behind \textbf{advent}.

\textbf{\sk{बुध्}} is PIE *\emph{bheudh-}, \textquotedblleft{}to be awake, to be aware.\textquotedblright{} Its most famous derivative is \textbf{\sk{बुद्ध}} (\emph{buddha}), \textquotedblleft{}the awakened one\textquotedblright{} — a title, not a name. Greek made it \emph{punthanomai}, \textquotedblleft{}to learn by asking\textquotedblright{}; Russian \emph{budit} still means \textquotedblleft{}to wake.\textquotedblright{}

English got it too, worn thin: \textbf{forbid} and \textbf{bode} (as in \emph{foreboding}) descend from Old English \emph{bēodan}, \textquotedblleft{}to announce, to offer.\textquotedblright{} Modern \emph{bid} mixes that verb with a second, unrelated one — so keep to \emph{forbid} and \emph{bode} if you want the clean cousin.
\end{cousinweb}

\subsection*{Your turn}
Say these aloud:

\begin{itemize}
\raggedright
  \item \textquotedblleft{}gacchati … avagacchati\textquotedblright{} — goes, understands
  \item the prefix and its sense — \emph{ava}, down
  \item the other verb and its root — \emph{budhyate}, \emph{budh}
  \item what \sk{बुद्ध} means — the awakened one
\end{itemize}

\begin{tcolorbox}[breakable,colback=teal!4,colframe=teal!35!black,title={Before you move on}]
Build \textquotedblleft{}understands\textquotedblright{} out of pieces you already had. (\textbf{\sk{अव}} \textquotedblleft{}down\textquotedblright{} + \textbf{\sk{गच्छति}} \textquotedblleft{}goes\textquotedblright{}.) Name the three verbs that share that one stem. (\emph{Gacchati}, \emph{āgacchati}, \emph{avagacchati}.) What does \textbf{\sk{बुध्}} mean before it means \textquotedblleft{}understand,\textquotedblright{} and what title comes from it? (\textbf{To wake}; \textbf{\sk{बुद्ध}}, the awakened one.) Which two English words are its clean cousins? (\textbf{Forbid} and \textbf{bode}.)
\end{tcolorbox}
\section[paṭhati]{\sk{पठति} (paṭhati) — \textquotedblleft{}reads,\textquotedblright{} and where the word went afterwards}

\label{lesson:SA-C08-pathati}



\begin{quote}
Say \textquotedblleft{}it thinks\textquotedblright{} and \textquotedblleft{}it understands.\textquotedblright{} (\emph{Cintayati}, \emph{avagacchati}.) This next verb is the plainest build in the chapter and the richest afterlife.
\end{quote}

\subsection*{You'll want to know: \sk{पठति}}
\noindent
\begin{tabularx}{\linewidth}{@{}>{\raggedright\arraybackslash}X>{\raggedright\arraybackslash}X>{\raggedright\arraybackslash}X@{}}
\toprule
\textbf{Sanskrit} & \textbf{sound} & \textbf{meaning} \\
\midrule
\sk{पठति} & \emph{paṭhati} & he, she, or it \textbf{reads} \\
\sk{पाठ} & \emph{pāṭha} & a \textbf{lesson}, a recitation \\
\bottomrule
\end{tabularx}

The root is \textbf{\sk{पठ्}} (\emph{paṭh}), class 1 — root, vowel, ending, exactly like \emph{bhavati}. Its first meaning is not \textquotedblleft{}read\textquotedblright{} but \textbf{\textquotedblleft{}recite aloud, rehearse, study\textquotedblright{}}: reading was a thing done with the mouth, and Sanskrit named it that way.

\begin{sounds}
\textbf{\sk{प}} is \emph{pa}. \textbf{\sk{ठ}} is \emph{ṭha}, a \emph{t} made with the tongue curled back to the roof of the mouth \emph{and} a puff of breath after it. \textbf{\sk{ति}} is \emph{ti}, tongue forward. Hear the two \emph{t}-sounds apart: \emph{pa-ṭha-ti}.
\end{sounds}

\begin{cousinweb}
This is the verb four living languages still read with — Hindi and Urdu \emph{paṛhnā}, Bengali \emph{poṛa}, Marathi \emph{paḍhṇe}. The change is regular: that curled-back \textbf{\sk{ठ}} softens between vowels, and the old ending drops off. A word reshaped by centuries of ordinary speech is \textbf{\sk{तद्भव}} (\emph{tadbhava}), \textquotedblleft{}become from that.\textquotedblright{}

Now set \textbf{\sk{चिन्ता}} beside it. It sits in Hindi, Marathi and Bengali \textbf{unchanged} — lifted back out of Sanskrit by scholars rather than inherited through speech. That kind is \textbf{\sk{तत्सम}} (\emph{tatsama}), \textquotedblleft{}the same as that.\textquotedblright{} Every claim about a Sanskrit word's descendants is one of those two, and they are not the same claim.
\end{cousinweb}

\begin{grammarlens}[title={Grammar Lens: where \sk{पठति} itself came from}]
The usual account runs back through \textbf{\sk{प्रथ्}} (\emph{prath}), \textquotedblleft{}to spread, to become widely known\textquotedblright{} — reciting as broadcasting — and behind that a root that also gives \textbf{\sk{पृथिवी}} (\emph{pṛthivī}), \textquotedblleft{}the earth, the wide one.\textquotedblright{}

Some hold instead that \emph{paṭhati} is a later, spoken form that scholars dressed up as Sanskrit. Keep the label on it, exactly as you did with \emph{punch}: a common account, not a settled one.
\end{grammarlens}

\subsection*{Your turn}
Say these aloud:

\begin{itemize}
\raggedright
  \item \textquotedblleft{}paṭhati\textquotedblright{} — it reads, it recites
  \item a worn-down grandchild — \emph{paṛhnā}
  \item a word carried whole — \emph{cintā}
\end{itemize}

\begin{tcolorbox}[breakable,colback=teal!4,colframe=teal!35!black,title={Before you move on}]
What did \emph{paṭhati} mean before \textquotedblleft{}read,\textquotedblright{} and which class builds it? (\textbf{Recite aloud}; \textbf{class 1}.) Name the two ways a Sanskrit word reaches a modern language. (\textbf{\sk{तद्भव}}, worn down through speech; \textbf{\sk{तत्सम}}, carried over whole.) Which is \emph{paṛhnā}, which is \emph{cintā}, and what label belongs on \emph{paṭhati}'s own ancestry? (\emph{Tadbhava}; \emph{tatsama}; \textbf{a common account, not settled} — the caution you kept on \emph{punch}.)
\end{tcolorbox}
\section[likhati]{\sk{लिखति} (likhati) — \textquotedblleft{}writes,\textquotedblright{} and the scratch it started as}

\label{lesson:SA-C08-likhati}



\begin{quote}
Say \textquotedblleft{}it reads.\textquotedblright{} (\emph{Paṭhati}.) Name a language that still reads with that word. (Hindi \emph{paṛhnā}, Bengali \emph{poṛa}, Marathi \emph{paḍhṇe} — any one.) Now the other half of literacy.
\end{quote}

\subsection*{You'll want to know: \sk{लिखति}}
\noindent
\begin{tabularx}{\linewidth}{@{}>{\raggedright\arraybackslash}X>{\raggedright\arraybackslash}X>{\raggedright\arraybackslash}X@{}}
\toprule
\textbf{Sanskrit} & \textbf{sound} & \textbf{meaning} \\
\midrule
\sk{लिखति} & \emph{likhati} & he, she, or it \textbf{writes} \\
\sk{लेखनी} & \emph{lekhanī} & a \textbf{pen} \\
\bottomrule
\end{tabularx}

The root is \textbf{\sk{लिख्}} (\emph{likh}), class 6. Class 6 adds the same \emph{-a-} class 1 does but leaves the root untouched: \emph{likh} + \emph{a} + \emph{ti}. Two classes, one vowel, and the difference is only whether the root gets strengthened first.

\begin{sounds}
\textbf{\sk{लि}} is \emph{li}. \textbf{\sk{ख}} is \emph{kha}, a \emph{k} with a puff of breath after it — the letter that opens \emph{khādati}. \textbf{\sk{ति}} is \emph{ti}: \emph{li-kha-ti}.
\end{sounds}

\begin{cousinweb}
\textbf{\sk{लिख्}} does not mean \textquotedblleft{}write\textquotedblright{} first. It means \textbf{to scratch, to scrape, to furrow}, and \textquotedblleft{}write\textquotedblright{} second — because writing \emph{was} scratching, into palm leaf and birch bark with a stylus. The older sense stays visible in \textbf{\sk{रेखा}} (\emph{rekhā}), \textquotedblleft{}a line, a scratch,\textquotedblright{} the same root wearing an \emph{r}.

Outside Sanskrit the trail is thin but real: Greek \emph{ereikō}, \textquotedblleft{}to tear,\textquotedblright{} and Lithuanian \emph{riekti}, \textquotedblleft{}to slice.\textquotedblright{} No everyday English word descends from it — that is two roots in this chapter that left English nothing.

Eastward, a different story: Hindi and Urdu \emph{likhnā}, Bengali \emph{lekhā}, Marathi \emph{lihiṇe} — worn down, \emph{tadbhava}, with Marathi's \emph{kh} softened to \emph{h}.
\end{cousinweb}

\begin{grammarlens}[title={What you've built: four verbs, four ways to build a stem}]
\emph{Cintayati}, \emph{avagacchati}, \emph{paṭhati}, \emph{likhati}. Thinking plants \textbf{-\sk{अय}-}; understanding is \textquotedblleft{}goes\textquotedblright{} with \textbf{\sk{अव}} in front, or \textbf{\sk{बुध्यते}}, waking; reading and writing take a plain \emph{-a-}, class 1 and class 6.
\end{grammarlens}

\subsection*{Your turn}
Say these aloud:

\begin{itemize}
\raggedright
  \item \textquotedblleft{}likhati\textquotedblright{} — it writes
  \item what the root meant first — to scratch
  \item the chapter's four — \emph{cintayati, avagacchati, paṭhati, likhati}
\end{itemize}

\begin{tcolorbox}[breakable,colback=teal!4,colframe=teal!35!black,title={Before you move on}]
What did \textbf{\sk{लिख्}} mean before \textquotedblleft{}write,\textquotedblright{} and what are the chapter's four verbs? (\textbf{To scratch}; \emph{cintayati, avagacchati, paṭhati, likhati}.) Which plants \textbf{-\sk{अय}-}, and which is a prefix on a verb you already had? (\emph{Cintayati}; \emph{avagacchati}.) Which verb of understanding means \textquotedblleft{}wakes,\textquotedblright{} and what title comes from its root? (\textbf{\sk{बुध्यते}}; \textbf{\sk{बुद्ध}}.) Is Bengali \emph{lekhā} tatsama or tadbhava, and which two word-histories here needed a label rather than a claim? (\textbf{Tadbhava}; \emph{\textbf{paṭhati}}'s ancestry and \emph{\textbf{likh}}'s English kin — the \emph{punch} caution, twice.)
\end{tcolorbox}
\section[ma]{\sk{म} — one character, met inside words you already say}

\label{lesson:SA-S01-letter-ma}



\begin{quote}
A horizontal head-line (shirorekha) runs across the top; letters hang beneath it like laundry on a line.

One character this time. Just one — and you have been saying it for pages without knowing which mark on the page it was.
\end{quote}

\subsection*{Script you'll notice: \sk{म}}
\textbf{\sk{म}} — \emph{ma}.

It is a \textbf{consonant}, and in this script a consonant is never bare: it comes with an \emph{a} already in it. So it is not \emph{m}, it is \textbf{ma}.

What it is made of:

\begin{itemize}
\raggedright
  \item a lower loop
  \item an upper loop
  \item right spine
  \item top bar
\end{itemize}

You already say these, and every one of them has \sk{म} somewhere inside it:

\begin{itemize}
\raggedright
  \item \textbf{\sk{नमस्ते}} \emph{namaste} — hello (lit. \textquotedblleft{}a bow to you\textquotedblright{}) (namaste)
  \item \textbf{\sk{नमस्कारः}} \emph{namaskāraḥ} — hello (lit. \textquotedblleft{}the making of a bow\textquotedblright{}) (namaskāraḥ)
  \item \textbf{\sk{स्वागतम्}} \emph{svāgatam} — welcome (lit. \textquotedblleft{}well come\textquotedblright{}) (svāgatam)
  \item \textbf{\sk{आम्} / \sk{न}} \emph{ām / na} — yes / no (ām / na)
\end{itemize}

\subsection*{Writing: \sk{म} — copy what you see}
Put your pen on \sk{म} and follow its line. Copy the shape you can see — slowly, and larger than it is printed.

\begin{quote}
This book does not yet tell you \textbf{where to start the character or which way to travel}. That is a real thing, taught with real variation from school to school, and it is not written down here until it can be written down with a source. Copying what is in front of you needs no such source, and it is how the shape gets into your hand in the meantime.
\end{quote}

\subsection*{Your turn}
\begin{itemize}
\raggedright
  \item \emph{Look:} at these words, and find \sk{म} in the ones that have it
\end{itemize}

\begin{quote}
\sk{नमस्ते}  ·  \sk{नमस्कारः}  ·  \sk{अस्ति}
\end{quote}

\begin{itemize}
\raggedright
  \item \emph{Trace it:} \sk{म} three times, saying \emph{ma} as you finish each one
  \item \emph{Look:} back at any page of this chapter and find \sk{म} once more
\end{itemize}

\begin{tcolorbox}[breakable,colback=teal!4,colframe=teal!35!black,title={Before you move on}]
Which character is this — \sk{म}? What sound does it carry? (\emph{\textbf{ma}}.) Name one word you already say that contains it.
\end{tcolorbox}
\section[ta]{\sk{त} — one character, met inside words you already say}

\label{lesson:SA-S107-letter-ta}



\begin{quote}
A horizontal head-line (shirorekha) runs across the top; letters hang beneath it like laundry on a line.

One character this time. Just one — and you have been saying it for pages without knowing which mark on the page it was.
\end{quote}

\subsection*{Script you'll notice: \sk{त}}
\textbf{\sk{त}} — \emph{ta}.

It is a \textbf{consonant}, and in this script a consonant is never bare: it comes with an \emph{a} already in it. So it is not \emph{t}, it is \textbf{ta}.

What it is made of:

\begin{itemize}
\raggedright
  \item a right-to-left upper shoulder curving down to an open lower tip
  \item a top-to-bottom right stem
  \item the top shirorekhā
\end{itemize}

You already say these, and every one of them has \sk{त} somewhere inside it:

\begin{itemize}
\raggedright
  \item \textbf{\sk{नमस्ते}} \emph{namaste} — hello (lit. \textquotedblleft{}a bow to you\textquotedblright{}) (namaste)
  \item \textbf{\sk{स्वागतम्}} \emph{svāgatam} — welcome (lit. \textquotedblleft{}well come\textquotedblright{}) (svāgatam)
  \item \textbf{\sk{अस्ति}} \emph{asti} — is
  \item \textbf{\sk{भवान्} / \sk{त्वम्}} \emph{bhavān / tvam} — you (respectful / familiar)
\end{itemize}

\subsection*{Writing: \sk{त}}
\begin{itemize}
\raggedright
  \item \textbf{1.} start at the body's upper-right junction, sweep left across the shoulder, then curve down around the left side and finish down-right at the open lower tip without lifting
  \item \textbf{2.} lift and draw the right stem top-to-bottom
  \item \textbf{3.} lift and draw the top shirorekhā left-to-right
\end{itemize}

\textbf{Pen lifts: 2.} The pen comes up 2 times and no more.

\begin{quote}
Verified three-stroke teaching form; everyday handwriting may join or simplify the body differently.
\end{quote}

\begin{quote}
This is one attested teaching order and not a national standard — handwriting here is taught with school-to-school variation. Source: Opiaterein, ‘Deva-\sk{त}-order.gif’, strokes 1–3, Wikimedia Commons, 10 May 2009.
\end{quote}

\subsection*{Your turn}
\begin{itemize}
\raggedright
  \item \emph{Look:} at these words, and find \sk{त} in the ones that have it
\end{itemize}

\begin{quote}
\sk{नमस्ते}  ·  \sk{स्वागतम्}  ·  \sk{नमस्कारः}
\end{quote}

\begin{itemize}
\raggedright
  \item \emph{Trace it:} \sk{त} three times, saying \emph{ta} as you finish each one
  \item \emph{Look:} back at any page of this chapter and find \sk{त} once more
\end{itemize}

\begin{tcolorbox}[breakable,colback=teal!4,colframe=teal!35!black,title={Before you move on}]
Which character is this — \sk{त}? What sound does it carry? (\emph{\textbf{ta}}.) Name one word you already say that contains it.
\end{tcolorbox}
